% !TeX root = ../papers/nc-p3.tex
\section{片側反転の最後の例外への全流入}
\label{sec:mixed-flip-fiber}
\begin{theorem}\label{thm:mixed-flip-fiber}
$w$ を $s=a=\iota$ の片側反転窓とする。全基底群の半安定商で $[w]$ の上にある
窓は全て標準Koszul左右AS正則である。同型類は以下の七座標の集合を以下の作用で割ったものに等しい。
座標を $b=(A,C,B,E,F,D,G)$ とし、この順の二次元トーラス指標を
\[
 (-1,0),\ (-1,2),\ (0,-1),\ (0,2),\ (1,1),\ (2,-1),\ (2,0)
\]
とする。許す台は $ABC$, $AC EF$, $CDEFG$ の部分集合、およびそれらの
反射 $A\leftrightarrow B,C\leftrightarrow D,E\leftrightarrow G$ である。
トーラス作用のほかに二作用 $Q,J$ で割る。
$Q$ はこの順に $(-1,-1,-1,1,1,-1,1)$ 倍、
$J$ は $(\iota,\iota,\iota,-\iota,1,\iota,\iota)$ 倍した後に上記の反射を行う。
一移動 $S$ は $(1,-1,-\iota,-1,-\iota,-\iota,1)$ 倍である。
最小周期は $Sb$ が $b$ と同値なら1、他は2である。
同じ五空間内外の同型・この周期条件はいずれも次の有限判定でよい。
$b'$ を $b,Qb,Jb,QJb$ のいずれかと比較し、台が一致し、
台上の指標行列の整数核の基底の各ベクトル $m$ について $\prod_j(b'_j/b_j^{*})^{m_j}=1$
となるか調べる。ここで $b^{*}$ は選んだ比較元である。
固定軌道の半安定な特殊化はこの同値を許し、台の指標が張る錐の任意の面の外の座標を零にしたものに尽きる。
\end{theorem}
\begin{proof}[五空間の構成と被覆]
前節の文字順で実際の四テンソルを与える再帰を記す。基底係数は $\mathbb Q(\iota)$ に置き、
共役転置による直交補をこの体で選んだ後 $k$ へ拡大する。
アフィン基底軌道の接空間を $O$、二切断の全接核を $K$ とする。
周囲の軌道横断面を $w+O^\perp$ に取り、$K\cap O^\perp$ の各一次元指標で
ガウス消去の自由成分を1にしたベクトル $b_j^0$ を取る。$O$ はスカラー方向も含む。
横断成分への直交射影を七座標と重み零の一座標に使う。
この選択は下記の位相付き単項行列と巡回に同変であり、トーラスについては重み空間が直交する。
周囲の作用写像の接は軌道に沿った因子を取れば同型であり、
閉軌道のLuna横断面定理で飽和開近傍に縮めると商からの基底変換として
この横断面を使える。基点へ縮む全軌道はその開近傍に入る。
安定群の恒等成分はトーラスなので、流入点はある余指標で正となる接座標だけを使う。
七つの内積の符号を平面で並べると極大の正台は上記五つで、
余指標は $(-3,-1),(-1,-3),(-1,2),(2,-1),(1,1)$ に取れる。
元の余指標での流入弧では周囲の座標が正の接座標の形式級数で決まり、
各項の指標はその非負和である。従って台を含む極大欄の余指標でも流入する。

各欄の座標を不定元とし、補正テンソル $D$ の単項式係数を総次数順に求める。
一次は $b_j^0$ とし、高次は $K^\perp$ のみに取る。各切断で $M$ を背景行列、
$L,R$ を左右核基底、$g$ を六つの独立行・列の交点上のピボット逆とする。
$MgM=M$ であり、六階数のSchur式は
$L(1+Dg)^{-1}DR=0$ である（ここでは $D$ は該当する切断を表す）。
指数 $m$ に対して $B_m=D_m-\sum_{u+v=m}D_u g B_v$ とすれば、
二切断の $LB_mR=0$ を同時に解けばよい。
最初の一次方程式 $D\mapsto (LDR)$ の像上で、$K^\perp$ への逆は
独立な行を $H$ として $H^*(HH^*)^{-1}$ で与える。
全ての行に戻して残差の消滅を確認する。各 $M,g$ は行列の源と先の
指標分解を保つ（独立行・列を選ぶ各ブロックは斉次）。
正の余重みしか使わないので $Dg$ は冪零であり、必要な単項式は各周囲語の指標に達するものだけで尽きる。
従ってこの再帰が通れば多項式による大域的な正規表示である。

\path{computations/verify_mixed_flip_fiber.py} は五つ全てで上の逆を取って
200行の残差が係数ごとに零になることを $\mathbb Q(\iota)$ 上で検算する。
この多項式表示の接射影は七つの入力そのものである。重なる台で高次補正は同じ再帰で一意である。
逆に、基点への流入弧の横断表示ではこの射影で全ての形式係数が一意に決まり、
従ってこの多項式表示が台を覆う。収縮の原点で六階数なので軌道上でも六階数、
命題\ref{prop:unrestricted-good-reduction}から全次数で正則となる。
\end{proof}
\begin{proof}[全安定群、一移動と軌道境界]
恒等成分の二文字・一文字・一文字の指標分解は内在的である。
従って二平面分解を保ち、二つの一文字を置換できる。
定理\ref{thm:split-spectral-line-classification}の交代成分 $02|13$ の二本の可分解直線と
第一の四対応により、二重文字の側も各位置で単項行列でなければならない。
ゆえに各位置各色の二文字の入替えの全 $2^8$ 通りを試せば十分である。
残る対角倍率は24個の支持語 $q$ と16個の位置文字 $(i,j)$ の行列
$N_{q,ij}=[q_i=j]$ で決まる。$N$ の非零Smith因子は1が十個と2が一個であり、
有限対角核はちょうど $Q$ で生成される（位置スカラーの積1は全空間に自明）。
上記検算では支持の一致後、係数比を四乗根で測り、Smith左変換で像への所属を判定し、
右変換で対角倍率を一組復元する。入替えで残るのは恒等と全位置の第一色のビット反転だけで、
後者の持上げが $J$ である。巡回したテンソルから同じ方法で一つの移送を選ぶと $S$ を得る。
全ての持上げは絶対値1の位相付き置換であり、スカラー規格化した背景を固定し、
横断接射影に上記の通り作用する。これで高次の座標変更も不要である。
$J^2$ はトーラスと $Q$ が生成する対角群に入り、$Q$ は反射不変である。
$S^2$ はトーラスの $(-1,1)$ と $Q$ の積に等しい。
代数閉体上で対角トーラスの像は指標の全整数関係で定まるので有限判定を得る。
横断面の商ファイバーと全基底群の商ファイバーの軌道は対応し、
その特殊化順序も一致する。対角トーラス軌道閉包の面による記述を有限個の移送と合わせれば最後の主張となる。
\end{proof}
